\chapter{Einführung die Turtle-Library und Schleifen}\label{ch:K01_Intro}

Um mit einem Computer zu kommunizieren, benötigt man eine Sprache, die er versteht – sogenannte Programmiersprachen. Durch das Schreiben von Programmen teilen wir dem Computer mit, welche Aufgaben er ausführen soll. In diesem Skript verwenden wir die (neueste Version) der Programmiersprache \textit{Python}.

Ähnlich wie natürliche Sprachen bestehen Programmiersprachen aus bestimmten Begriffen mit festgelegter Bedeutung. Begriffe, die Anweisungen zur Ausführung von Befehlen darstellen, werden Befehlswörter oder kurz Befehle genannt.

Ein Programm ist eine Sammlung von Befehlen in einer Programmiersprache, die eine bestimmte Aufgabe beschreiben. Grundsätzlich kann man sich ein Programm als eine Abfolge von Anweisungen vorstellen, die der Computer Schritt für Schritt verarbeitet.

Das Ziel des Programmierens besteht darin, Abläufe zu automatisieren. Dabei überträgt man die Ausführung einer Aufgabe vollständig an den Computer. Deshalb muss ein Programm so gestaltet sein, dass es eine eindeutige Bedeutung hat und exakt beschreibt, was getan werden soll.

In diesem Kapitel werden Sie Ihre ersten Programme schreiben und das Konzept der Schleifen kennenlernen. Schleifen ermöglichen es, wiederkehrende Aufgaben mehrfach automatisch ausführen zu lassen.

\begin{myremark}
	Unter dem Link
	\begin{center}
		\url{https://docs.python.org/3/library/turtle.html#module-turtle}
	\end{center}
	finden Sie eine Übersicht aller Turtle-Befehle. Wir empfehlen Ihnen für diesen Link ein Lesezeichen (Bookmark) in Ihrem Webbrowser zu erstellen.
	Einige dieser Befehle werden wir im Folgenden gemeinsam genauer kennenlernen.
\end{myremark}
\lstinputlisting[language=Python,caption=\texttt{strich.py},label=listing:strich]{GrundlagenInfo/00_Programmieren/Skript/Code/K01_Intro/strich.py}
\lstinputlisting[language=Python,caption=\texttt{dreieck.py},label=listing:dreieck]{GrundlagenInfo/00_Programmieren/Skript/Code/K01_Intro/dreieck.py}
